\section{The complete relative strong first-energy space}
\label{sec:r11-strong-first-energy-completion}

Retain precisely the strict analytic box and extension hypotheses
(TC1)--(TC2), the arbitrary original separately closed symmetric
middle-positive background $S$, and $q=n-p\geq1$. Write
$d=q+1$, $\alpha=2^{-q}$, and $\mathcal E^1=
\mathcal E^1_{S,\partial}(V)$. The actual classical class
$\mathcal E^1$ is not enlarged by definition. We instead construct
a distinct space of actual Borel functions modulo local ordinary-open
outer $S$-capacity. Its defining rows have no prescribed common
actual lower barrier. We do not assert that its members admit
ambient classical representatives or compact negative sublevel
closures. Their boundary condition is the strong approximation
condition printed below.

The frozen bounded predecessors are the R2 capacitary completion,
R3 bounded common-collar strong topology, and R4 completed
variational retraction. The new inputs are the unbounded global
first-energy calculus, the tester estimate (EV1), and the analytic
bounded completion in Lemma~\ref{lem:r11-analytic-bounded-completion}.
The present construction does not reinterpret the frozen linear
Hilbert completion as a nonlinear classical potential space.

\subsection{Uniform tails and the boundary of an energy-Cauchy row}

Call an actual row $(u_j)\subset\mathcal E^1$ strongly Cauchy if
\[
               \lim_{N\to\infty}\sup_{j,l\geq N}
                    \mathscr I_S(u_j,u_l)=0.
 \tag{IC1}
\]
No triangle inequality for $\mathscr I_S$ is assumed.

\begin{lemma}[Uniform deep tails and global weighted tightness]
\label{lem:r11-strong-cauchy-uniform-tails}
Every strongly Cauchy actual row has bounded first energies and
\[
 \begin{split}
 \lim_{R\to\infty}\sup_j
     \int_{\{u_j<-R\}}(-u_j)\,d\mu_{u_j}&=0,\\
 \lim_{R\to\infty}\sup_j
     \int_R^\infty t^q C_S(\{u_j<-t\};V)\,dt&=0.
 \end{split}
 \tag{IC2}
\]
For finitely many strongly Cauchy actual rows $u_{i,j}$,
$0\leq i\leq q$, the finite positive Radon measures
\[
            (-u_{0,j})M(u_{1,j},\ldots,u_{q,j})
 \tag{IC3}
\]
are uniformly tight on $V$. The same conclusions hold if the
indices of the finitely many rows tend independently to infinity.
\end{lemma}

\noindent\textbf{Status: \texttt{NEW\_PROVED}}. The first assertion
extends (EV19) from convergence to a specified actual member to
arbitrary strongly Cauchy rows; the boundary tightness is proved
without a common lower barrier.

\begin{proof}
The energy-root estimate (EV9) bounds the energies by some $M$.
The case $M=0$ is immediate from the constant-one mixed Holder
inequality; otherwise enlarge $M$ to be positive. On
$\{-u_j>R\}$ split according to whether
$-u_N\leq(-u_j)/2$. The calculation used for (EV20) gives
\[
 \int_{\{u_j<-R\}}(-u_j)\,d\mu_{u_j}
 \leq 2\int|u_j-u_N|\,d\mu_{u_j}
       +4\int|u_N-u_N^{(R/4)}|\,d\mu_{u_j},
 \qquad u^{(L)}=\max(u,-L).
 \tag{IC4}
\]
By (EV1), the first integral is at most
$C_{q,M}\mathscr I_S(u_j,u_N)^\alpha$.
For fixed $N$, the second is at most
$C_{q,M}\mathscr I_S(u_N,u_N^{(R/4)})^\alpha$, uniformly
in $j$. The latter discrepancy tends to zero by (EV15).
Choose $N$ large using (IC1), then let $R$ grow, and finally
make the Cauchy error tend to zero. The finitely many indices
$j<N$ have vanishing tails because their first energies are
finite. This proves the first line of (IC2). The tail form of
(EV21),
\[
 \int_R^\infty t^q C_S(\{u_j<-t\};V)\,dt
       \leq 2^{q+1}\int(-u_j-R/2)^+\,d\mu_{u_j},
 \tag{IC5}
\]
proves the second. For a finite sum of rows, the sublevel union
bound gives the same uniform tail conclusion, with the factor
$(\#\text{slots})^{q+1}$ and a fixed rescaling of $R$.

For (IC3), all slot energies are uniformly bounded. Write
$M_j=M(u_{1,j},\ldots,u_{q,j})$. Positivity of the expansion
of $\mu_{\sum_{i=1}^q u_{i,j}}$ and (EV1) give
\[
             \int|u_{0,j}-u_{0,N}|\,dM_j
                  \leq C\mathscr I_S(u_{0,j},u_{0,N})^\alpha,
 \tag{IC6}
\]
where $C$ depends on the uniform energies and $q$, but not on
the derivative-row indices. Fix a sufficiently late $N$.
For $\epsilon>0$, the set
$K_\epsilon=\overline{\{u_{0,N}<-\epsilon\}}$ is compact.
On its complement $u_{0,N}=u_{0,N}^{(\epsilon)}$, and hence
\[
 \begin{split}
 \int_{V\setminus K_\epsilon}(-u_{0,j})\,dM_j
 &\leq \int|u_{0,j}-u_{0,N}|\,dM_j
          +\int(-u_{0,N}^{(\epsilon)})\,dM_j\\
 &\leq C\mathscr I_S(u_{0,j},u_{0,N})^\alpha
       +e_1(u_{0,N}^{(\epsilon)})^{1/d}
                          \prod_{i=1}^q e_1(u_{i,j})^{1/d}.
 \end{split}
 \tag{IC7}
\]
The last line is the global constant-one Holder inequality.
The small-value energy vanishes:
\[
 e_1(u_{0,N}^{(\epsilon)})
   \leq C_q\int_0^\epsilon t^q
                    C_S(\{u_{0,N}<-t\};V)\,dt\longrightarrow0,
 \tag{IC8}
\]
by the power-energy/capacity equivalence. First make the
Cauchy error small, then choose $\epsilon$ small. The
finitely many omitted weights $u_{0,j}$, with derivative
rows still varying, are handled by the same argument with
zero first term: outside their own compact negative sublevel,
the Holder bound with $u_{0,j}^{(\epsilon)}$ tends uniformly
to zero. A finite union of these compacts proves uniform
tightness. This observation also proves the assertion for
independent row indices, without taking a diagonal assumption
on any derivative slot.
\end{proof}

\subsection{The Borel class and its intrinsic original currents}

\begin{definition}[Strong relative first-energy completion]
\label{def:r11-strong-first-energy-completion}
Let $\widehat{\mathcal E}^1_{\rm str}(V)$ consist of actual
Borel functions $u:V\to[-\infty,0]$ which are limits of an
actual strongly Cauchy row in every local truncated outer
$S$-capacity. Equality means local outer-capacity equality.
The capacities may equivalently be taken on all compact
cores in the fixed domain $V$, by (AL9). A choice of actual
defining row is not part of the element of this quotient.
\end{definition}

\begin{proposition}[Borel limits of all strongly Cauchy actual rows]
\label{prop:r11-strong-cauchy-borel-limit}
Every row satisfying (IC1) has a limit in
Definition~\ref{def:r11-strong-first-energy-completion}, unique
modulo local outer-capacity zero. Every member of that space is
finite and continuous locally outside ordinary open sets of
arbitrarily small outer capacity. Its pole set is locally
outer-capacity zero.
\end{proposition}

\noindent\textbf{Status: \texttt{NEW\_PROVED}} for unrestricted
strongly Cauchy first-energy rows, with no common actual barrier.

\begin{proof}
Uniform energy and (EV8) make each fixed bounded truncation
of the row Cauchy in every local outer capacity. Apply the
countable compact-core and integer-level diagonal construction
from the bounded completion. A subsequence has compatible
pointwise limits of all bounded truncations off a locally
outer-null set. Its actual pointwise limsup defines a Borel
$u\leq0$ on all of $V$, and the Cauchy property gives the
required capacity convergence for the full row.

There is a uniform deep sublevel estimate. The no-loss
sublevel inequality and the first-energy bound imply
\[
            C_S(\{u_j<-R\};V)\leq C_q M R^{-d}.
 \tag{IC9}
\]
These sets are ordinary open, so the same bound holds for
their outer capacities. On a fixed compact core, the set
$\{u<-2R\}$ is contained in $\{u_j<-R\}$ together with
a positive gap of the fixed truncations at $3R$.
Pass $j\to\infty$ in the outer-capacity subadditivity bound.
It follows that the pole of $u$ has outer capacity zero on
that core. Outside an ordinary open cover of a sufficiently
deep sublevel, $u$ agrees with one bounded truncation.
Its bounded-completion quasicontinuity then proves the stated
finite local quasicontinuity. The level-set triangle inequality
proves uniqueness. Neither this argument nor its conclusion
requires compact negative sublevel closures for the limit.
\end{proof}

\begin{theorem}[Canonical critical calculus in the strong completion]
\label{thm:r11-strong-completed-calculus}
For a finite tuple $\mathbf u$ in
$\widehat{\mathcal E}^1_{\rm str}(V)$, with actual strongly
Cauchy defining rows $\mathbf u_j$, there are intrinsic
coefficient-Radon original currents
\[
 \widehat N_\gamma(\mathbf u),\qquad
 \widehat{\mathcal J}'_{\gamma,i},\quad
 \widehat{\mathcal J}''_{\gamma,i},\qquad
 \widehat G^\gamma_{ih}.
 \tag{IC10}
\]
Curvatures have $k=|\gamma|\leq q$; first-order and energy
currents have $k\leq q-1$. All energy degrees are included,
since $k+2\leq q+1=d$. Curvatures are separately closed,
symmetric and middle-positive; energy matrices are positive.
First-order and energy currents are not asserted to be closed.

For any continuous $F$ on the nonpositive orthant satisfying
$|F(x)|\leq A(1-\sum_i x_i)^a$, the actual multiplier limit
\[
                  F(\mathbf u_j)P_j
                         \longrightarrow F(\mathbf u)\widehat P
       \qquad(a+b\leq d)
 \tag{IC11}
\]
holds on every original compact smooth complementary test and
every actual moment refinement. Here the degree budget $b$ is
$k$, $k+1$, or $k+2$ for curvature, first-order, or energy.
All these multipliers are coefficient-integrable. The currents
and weighted limits are independent of defining rows and of
outer-capacity-null representatives, and they annihilate the
joint pole set. They satisfy all original recursive identities,
including
\[
 \begin{split}
 D(u_i\widehat N_\gamma)&=\widehat N_{\gamma+e_i},\\
 d'(u_i\widehat N_\gamma)&=\widehat{\mathcal J}'_{\gamma,i},\qquad
 d''(u_i\widehat N_\gamma)=\widehat{\mathcal J}''_{\gamma,i}
                         \qquad(k\leq q-1),
 \end{split}
 \tag{IC12}
\]
and the actual three-term energy identity (RC13).

For $q+1$ completed slots, the global mixed energy is finite and
\[
 \mathcal I_S(u_{0,j},\ldots,u_{q,j})
      \longrightarrow
 \widehat B(u_0,\ldots,u_q)
       :=\int_V(-u_0)\widehat N(u_1,\ldots,u_q).
 \tag{IC13}
\]
The corresponding positive weighted measures converge against
every bounded original smooth scalar test, including those
without compact support. The assertion also permits independent
indices tending to infinity in the defining rows.
\end{theorem}

\noindent\textbf{Status: \texttt{NEW\_PROVED}} for the strong
completion without a prescribed actual lower barrier.

\begin{proof}
For a finite tuple define
\[
 \Phi(R)=\sup_j\int_R^\infty t^q
          C_S\!\left(\left\{\sum_i u_{i,j}<-t\right\};V\right)dt.
 \tag{IC14}
\]
By Lemma~\ref{lem:r11-strong-cauchy-uniform-tails},
$\Phi(0)<\infty$ and $\Phi(R)\to0$; the same is true for
the supremum over independently large row indices, with the
finitely many earlier indices included if needed. At every
original compact test stage, the actual truncation estimates
(PG3), (FO8), and (PG12) give
\[
 \|\mu_{P_j-P_j^{(L)},\eta}\|_{\rm TV}
             \leq C_\eta L^{b-d}\Phi(L/4).
 \tag{IC15}
\]
The bounded analytic completion supplies intrinsic currents
$P^{(L)}$ of the completed clipped tuple. Passing $j$ to
infinity in the difference of two bounded stages shows that
$P^{(L)}$ are Cauchy in local coefficient total variation,
exactly as in (RC10). Their limit defines $\widehat P$.
Coefficient/refinement identities, positivity and separate
closedness pass on all original tests.

The limiting coefficients have the common modulus
\[
    \inf_{L\geq1}\{L^b s+L^{b-d}\Phi(L/4)\}
                              \longrightarrow0\quad(s\downarrow0).
 \tag{IC16}
\]
Its weak-limit proof keeps $L$ fixed until after ordinary-open
variation lower semicontinuity and Radon outer regularity.
Thus it controls both the defining and limiting coefficients,
including the critical case $b=d$. It annihilates every
locally outer-null set and, in particular, the joint pole set.
The bounded quasicontinuous multiplier argument therefore
passes every bounded continuous function of a fixed clipped
tuple.

For endpoint integrability apply that argument to
$\theta_R(z)\min((1+z)^a,H)$, where
$z=-\sum_i u_i$ and $\theta_R$ is zero below $R$ and one
above $2R$. For fixed $R,H$ this is a continuous function
of one sufficiently large clipped tuple; once a slot exceeds
that clipping level its value is the constant $H$.
When $a=0$, use $\theta_R$ itself. The actual row estimates
bound the integral by
$C R^{a+b-d}\Phi(R/8)$. Passing to the bounded multiplier
limit and then $H\uparrow\infty$ proves the same tail bound
for the limiting positive traces. Variation lower
semicontinuity handles signed first-order coefficients, and
positive diagonal controls handle cross energies.
Adding the bounded-value part gives actual coefficient
integrability. Clipping a general $F$ and then letting its
clipping level tend to infinity proves (IC11). Consequently
all multipliers in (IC12) and (RC13) are defined before
distributional differentiation, and the original row
identities pass to the limit.

For two different defining tuples with the same Borel
capacity limits, their bounded completed stages are identical.
The union of their finitely many row families has the same
vanishing-tail property. The preceding construction therefore
gives identical unbounded currents and actual multipliers.
This proves approximation independence without presupposing
that the two tuples are mutually strongly Cauchy.

Finally, (IC11) gives local convergence of (IC3).
The uniform global tightness in
Lemma~\ref{lem:r11-strong-cauchy-uniform-tails}, together
with the uniform finite mixed-energy bound (GE10), identifies
a finite global limiting measure and gives its total-mass
convergence. This is precisely (IC13). Tightness also passes
every bounded original scalar test. The independent-index
assertions follow from the same bounds, or by applying the
argument to every sequence of independently diverging indices.
\end{proof}

\subsection{Intrinsic discrepancy, completeness, and the exact scope}

Define on this Borel space
\[
 \begin{split}
 \widehat e(u)&=\widehat B(u^{[q+1]}),\\
 \widehat I(u,v)&=\widehat e(u)+\widehat e(v)
       -\widehat B(u,v^{[q]})-\widehat B(v,u^{[q]}).
 \end{split}
 \tag{IC17}
\]

\begin{lemma}[Uniform discrepancy continuity on energy balls]
\label{lem:r11-discrepancy-uniform-perturbation}
The completed space is closed under finite sums and nonnegative
scalar multiples. Its mixed energy is symmetric, positive and
multilinear on that cone, and satisfies the constant-one Holder
inequality. The formulas (EV1), (EV8), (EV9) and the sharp
midpoint inequality (SM8) hold in the completion, with their
actual Borel multipliers and local capacities. In particular,
$\widehat I\geq0$, and it vanishes exactly on equal relative
classes.

If $\widehat e(a),\widehat e(a'),\widehat e(b)\leq M$,
then
\[
 \big|\widehat I(a,b)-\widehat I(a',b)\big|
       \leq L_q M^{1-\alpha}\widehat I(a,a')^\alpha,
 \qquad
 L_q=2(q+1)2^{(q-1)\alpha}3^{d(1-\alpha)}.
 \tag{IC18}
\]
The assertion at $M=0$ means that both sides are zero.
\end{lemma}

\noindent\textbf{Status: \texttt{NEW\_PROVED}} for the uniform
discrepancy modulus; the global algebraic inequalities are
\texttt{DIRECT\_COROLLARY} extensions through (IC13).

\begin{proof}
For two strongly Cauchy defining rows, the sum rows have the
required Borel capacity limit. Expand the discrepancy of two
independently late sum rows into finitely many global mixed
energies. Equation (IC13) makes each term converge to the
corresponding completed mixed energy; the total limiting
discrepancy is zero. Thus the sum row is strongly Cauchy.
Positive scaling is immediate, and zero scaling means the
zero function even at poles. Passing the actual finite mixed
identities and (GE10) through (IC13) now proves all asserted
cone identities.

The actual inequalities (EV1), (EV9), and (SM8) pass by
global mixed-energy convergence. For (EV1), its left side is
the actual weight $|u-v|$, which has degree-one growth and is
bounded by $(-u)+(-v)$; (IC11) and the tightness of these two
critical positive mixed measures pass its global integral.
For (EV8), use defining rows, their convergent energies,
and the level-set triangle inequality for clipped gaps;
then remove deep sublevels using (IC9) and decrease the
intermediate gap threshold. This also proves that
$\widehat I=0$ gives equality in all local outer capacities.
The converse follows from intrinsic current and multiplier
independence.

It remains to prove (IC18). It suffices first to consider
actual potentials. Write $h=a-a'$. Telescoping the three
terms of $I(a,b)$ that depend on $a$, and using full mixed
symmetry, expresses their difference as $2(q+1)$ signed
integrals of $-h$ against positive mixed top products whose
slots lie among $a,a',b$. Every such product is bounded by
$\mu_\psi$, where $\psi=a+a'+b$, by the positive expansion.
Thus
\[
 |I(a,b)-I(a',b)|
                  \leq2(q+1)\int|a-a'|\,d\mu_\psi.
 \tag{IC19}
\]
Minkowski gives $e_1(\psi)\leq3^dM$ and
$e_1(a+a')\leq2^dM$. Apply (EV1) to (IC19) to obtain
(IC18). For completed potentials pass through actual
defining rows whose energies tend to the completed energies;
first allow $M+\epsilon$ and then let $\epsilon\downarrow0$.
The zero-energy case follows separately from Holder.
\end{proof}

\begin{theorem}[Completeness and the intrinsic strong topology]
\label{thm:r11-complete-strong-first-energy-space}
Every $\widehat I$-Cauchy sequence in
$\widehat{\mathcal E}^1_{\rm str}(V)$ converges to a unique
relative class in that space, with discrepancy tending to zero.
The actual classical class is strongly dense.

For members $u_j,u$ of this space, one has
\[
 \widehat I(u_j,u)\longrightarrow0
 \quad\Longleftrightarrow\quad
 \left\{
 \begin{array}{l}
 u_j\longrightarrow u\text{ in all local truncated outer capacities},\\
 \widehat e(u_j)\longrightarrow\widehat e(u).
 \end{array}\right.
 \tag{IC20}
\]
Strongly convergent finite tuples have every original local
weighted-current limit in (IC11), including all its critical
endpoints, and the global mixed-energy and weighted-measure
limits (IC13). These assertions do not claim local total
variation for arbitrary varying strongly convergent rows.

Every fixed actual lower-barrier completion
$\widehat{\mathcal E}^1_a$ of (RC2) embeds canonically in this
space, preserving all its original currents, actual multipliers
and mixed energies.
\end{theorem}

\noindent\textbf{Status: \texttt{NEW\_PROVED}} for the full relative
strong completion. This does not assert actual classical
completeness.

\begin{proof}
If $(v_l)$ is an actual defining row for a completed $u$, then
for fixed $l$, (IC13) gives
$\widehat I(v_l,u)=\lim_{r\to\infty}I(v_l,v_r)$.
Its Cauchy property therefore implies
$\widehat I(v_l,u)\to0$. This proves strong density.

Now let $(u_j)$ be completed and $\widehat I$-Cauchy.
The energy-root estimate bounds its energies. Choose actual
$v_j$ with $\widehat I(v_j,u_j)<2^{-j}$ and uniformly bounded
energies. Applying (IC18) in its two slots shows
\[
 I(v_j,v_l)\leq\widehat I(u_j,u_l)
                   +C\bigl(2^{-j\alpha}+2^{-l\alpha}\bigr).
 \tag{IC21}
\]
Hence $(v_j)$ is strongly Cauchy and has a completed Borel
limit $u$. Its defining property and strong density imply
$\widehat I(v_j,u)\to0$. A further use of (IC18) gives
$\widehat I(u_j,u)\to0$. Separation proves uniqueness.
This proves completeness using a proved uniform modulus,
rather than an unproved triangle inequality.

The forward implication of (IC20) follows from the completed
energy-root and capacity estimates. For the reverse, the
energies are uniformly bounded. Bounded completed stages and
the strict subcritical truncation bound for top unweighted
currents give local convergence of their top measures. More
explicitly, the stage error has order at most $C M/L$ for
degree $q<d$; first fix $L$, use bounded-completion capacity
continuity, and then let $L$ grow. The same uniform coefficient
modulus permits bounded quasicontinuous weight passage.
Consequently, by first truncating the weight and inserting a
compact nonnegative cutoff, then increasing both,
\[
 \liminf_j\widehat B(u_j,u^{[q]})\geq\widehat e(u),\qquad
 \liminf_j\widehat B(u,u_j^{[q]})\geq\widehat e(u).
 \tag{IC22}
\]
Holder gives upper bounds tending to $\widehat e(u)$ for
both cross energies. Expanding (IC17) proves the converse.
This argument uses no critical weighted continuity to prove
that very continuity.

For a strongly convergent completed tuple choose, at the $j$th
row, classical approximants with discrepancy at most $2^{-j}$,
and sufficiently accurate values on the first $j$ compact
test stages and bounded clipped weights. The defining-row
convergence allows these finitely many requirements
simultaneously. By (IC18), these actual rows are strongly
Cauchy, with the same Borel limits and energies. The preceding
construction and (IC11)--(IC13) give the asserted limits.
Equivalently, any failure on a fixed test would furnish a
subsequence and actual approximants contradicting those
already proved actual-row limits. This last argument also
covers bounded original scalar tests in the global statement.

Finally, classical defining rows of a fixed-barrier member
have joint critical mixed-energy convergence by (RC14), even
when two row indices tend independently. Expanding their
discrepancy therefore makes them strongly Cauchy. The Borel
representative is unchanged, and the intrinsic bounded stages
agree in the two constructions. Hence all local currents,
actual weighted multipliers and global mixed energies agree.
\end{proof}

\paragraph{The remaining source and representative gates.}
The completed space solves the strong Cauchy problem for the
relative first-energy geometry. It does not turn a capacity
class into an actual ambient classical psh function. It also
does not, by itself, extend an arbitrary datum integral
$u\mapsto\int u\,d\mu$ continuously to this Borel space.
Such a datum statement requires its own null-set and uniform
integrability hypotheses; finite energy-dual size alone is
not used as a substitute for them. The barrier-constrained
problem (BV7)--(BV8) retains its printed integrable barrier
and BT-diffuse datum assumptions. No unrestricted nonlinear
source equation, GQ, or low-real-degree CTO is inferred.
